\subsection{A unified quadratic and sphericalization continuation}
\label{sec:m5-quadratic-continuation}
The continuation of 20 October applies three related mechanisms: the
full induced-unit compression bound, exhaustive sphericalization followed
by dimension-field reduction, and simultaneous first and second central
twist moments.  Their roles are different.  Compression can exclude an
arbitrary categorification directly; a spherical obstruction requires a
separate argument covering every signed quotient; and a pseudo-unitary
center obstruction additionally requires the spherical dimension
character to belong to the Frobenius--Perron Galois orbit.

\begin{proposition}\label{prop:m5-direct-compression}
The rings $R7\_5\_4653$, $R7\_5\_5029$ and $R7\_5\_5040$
admit no complex fusion categorification.
\end{proposition}
\begin{proof}
For a real quadratic field $L$ of fundamental discriminant $q$, collect
the irreducible factors $(x^2-t_jx+n_j)^{m_j}$ of the regular Casimir
characteristic polynomial whose roots generate $L$.  Define
\[
 Q_L(R)=q\sum_j m_j\frac{t_j^2-4n_j}{n_j^2}.
\]
Corollary~\ref{cor:compression-boundary} gives the necessary condition
$Q_L(R)\leq4$ for arbitrary complex categorification.  In particular, it
does not require the original category to admit a pivotal structure.
The exact factors and rational values are
\begin{center}\small
\begin{tabular}{@{}l l r r@{}}
\toprule
Ring & Quadratic Casimir factor & $q$ & $Q_L(R)$\\
\midrule
$R7\_5\_4653$ & $x^2-3960x+7260$ & $165$ & $49$\\
$R7\_5\_5029$ & $x^2-1584x+6336$ & $88$ & $49/9$\\
$R7\_5\_5040$ & $x^2-1584x+6336$ & $88$ & $49/9$\\
\bottomrule
\end{tabular}
\end{center}
Each exceeds four.  Recall why the general bound differs from the
spherical bound $[K:L]Q_L(R)\leq1$: in sphericalization the unsigned
sector has Casimir $2B$, so its contribution is $Q_L(R)/4$.  The other
sector contributes a nonnegative squared norm.  The obstruction therefore
holds independently of any proposed choice of signed quotient.
\end{proof}

\begin{proposition}\label{prop:m5-signed-compression}
The ring $R7\_5\_4569$ admits no complex fusion categorification.
\end{proposition}
\begin{proof}
Its Casimir characteristic polynomial is
\[
 (x-14)(x-3)(x-7)^3(x^2-2940x+17640).
\]
For $L=\mathbb Q(\sqrt{15})$, of fundamental discriminant $60$, this
gives $Q_L(R)=81/49>1$.  Corollary~\ref{cor:quadratic-Casimir} excludes
a spherical categorification, but this inequality alone would not
exclude an arbitrary fusion category.

Apply the complete oriented sphericalization enumeration to the original
tensor.  All its simples are self-dual.  Choose the lift with positive
second indicator for each simple, with the actual tensor unit as unit.
The difference tensor $T$ has the same parity as $N$, satisfies
$-N_{ij}^k\leq T_{ij}^k\leq N_{ij}^k$, and obeys the oriented cyclic
identities.  The odd-square-channel condition from
\cite[Theorem~3.9]{TPE} imposes $T\equiv N\pmod4$ on the relevant
channels.  These constraints apply in sphericalization and do not
assume a pivotal structure on the source category.

Exhausting associativity on these finite domains gives exactly one
solution, namely $T=N$.  The replay visits $2285932869$ search nodes;
completion is regenerated, rather than inferred from a candidate list.
The multiplicity of an opposite lift is $(N-T)/2$, hence is zero in every
product.  The selected lifts form a spherical tensor subcategory, and
forgetting is a tensor equivalence to the original category.  This
contradicts its spherical obstruction.
\end{proof}

\begin{proposition}\label{prop:m5-quadratic-galois-pu}
Every complex fusion categorification of any of
\[
 R7\_5\_443,\quad R7\_5\_1214,\quad
 R7\_5\_3895,\quad R7\_5\_5039
\]
admits a spherical structure and has a pseudo-unitary Galois conjugate.
This is a structural reduction, not an existence assertion.
\end{proposition}
\begin{proof}
Complete oriented enumeration gives respectively $1,1,2,4$ signed
quotients.  For the first two rings the sole quotient is the original
tensor.  Both quotients of $R7\_5\_3895$, and two quotients of
$R7\_5\_5039$, are gauge equivalent to the original tensor.  Switching
the lift choices by that gauge gives a tensor section of sphericalization,
as in Proposition~\ref{prop:m5-signed-compression}.

For either remaining quotient of $R7\_5\_5039$, the unsigned and signed
Casimirs contain respectively
\[
 x^2-1440x+6400,\qquad x^2-616x+1760.
\]
Their roots generate $\mathbb Q(\sqrt5)$ and $\mathbb Q(\sqrt{11})$.
The two sectors of the sphericalization have twice these Casimir roots.
The algebras are commutative, so these are twice the corresponding formal
codegrees.  Every formal codegree of a spherical category belongs to its
dimension field: a constituent of the induced unit has dimension $D/f$,
and forgetting expresses this dimension as an integer combination of
source simple dimensions.  Thus the sphericalization's dimension field
must contain the degree-four compositum of these two quadratic fields.
But every character component of either signed quotient has degree at
most two.  Its spherical dimension character belongs to that signed
sector, giving a contradiction.  This excludes the nongauge quotients.

Finally the exact character decomposition of each original ring shows
that every character outside the Frobenius--Perron orbit either vanishes
on a simple object or takes unequal values on dual objects.  Neither is
a spherical dimension character.  Galois conjugation therefore makes
all simple dimensions equal to their Frobenius--Perron values, proving
the assertion.
\end{proof}

\subsection{The conductor-sixty two-moment obstruction}
\label{sec:m5-conductor-sixty}
\begin{theorem}\label{thm:m5-q60}
The rings $R7\_5\_443$ and $R7\_5\_1214$ admit no complex fusion
categorification.
\end{theorem}
By Proposition~\ref{prop:m5-quadratic-galois-pu}, it suffices to exclude
a pseudo-unitary categorification with its canonical spherical structure.
The following argument extends Section~\ref{sec:quadratic-center-trace}
to conductor sixty.  A relative-trace refinement avoids any assumption
that paired individual twists coincide.

Put $s=\sqrt{15}$ and $\eta=4+s$.  The source dimension vectors are
\[
\begin{array}{c|l}
443&(1,1,4+s,5+s,5+s,4+s,6+2s)\\
1214&(1,1,4+s,4+s,5+s,5+s,6+2s).
\end{array}
\]
Both global dimensions are $D=60\eta$, with $\eta\bar\eta=1$.
The quadratic Casimir factor is $x^2-480x+3600$, so $Q_{60}=1$.
The remaining formal codegrees are respectively
$(12,5,4,6,6)$ and $(5,12,12,4,4)$.

There is a cyclotomic automorphism that flips $s$ and fixes every other
source character.  To check the latter assertion, use the separating
matrix $\sum_i w_iN_i$, where $w=(1,3,11,37,101,307,997)$.
Its characteristic polynomials are
\begin{align*}
443:\quad &(x-683)(x-422)(x-230)(x-184)(x+298)\\[-2pt]
&\hspace{6mm}\cdot(x^2-15896x-26866796),\\
1214:\quad &(x-953)(x-24)(x+28)\\[-2pt]
&\hspace{6mm}\cdot(x^2-16436x-22501976)(x^2+712x-572).
\end{align*}
All roots are distinct.  The only other nonrational character field is
$\mathbb Q(\sqrt3)$ in the second ring, which is disjoint from
$\mathbb Q(\sqrt{15})$.  Formal-codegree rationality alone would not
suffice for this choice of automorphism.

Write $d_i=a_i+b_is$, $a(v)=\sum_i a_iv_i$ and $b(v)=\sum_i b_iv_i$.
The full character transport formula
\eqref{eq:quadratic-row-galois} gives the integral involution
\[
                  T(v)=v+(b_i)_i\bigl(a(v)-4b(v)\bigr).
\]
Every central restriction row has partner $T(v)$ or $-T(v)$, with one
common sign making all entries nonnegative.  A mixed-sign image is
impossible.  The Galois image $J$ of the central unit has restriction row
$(1,0,1,1,1,1,2)$ and dimension $\eta^2=31+8s$.

For $w=e_{\mathbf1}-e_J$, the known induced-unit coordinates give
\[
 \operatorname{proj}_{I(1)}Sw
 =(D^{-1}-\bar D^{-1})w=-\frac{w}{\sqrt{60}}.
\]
This saturates the compression bound.  The equality and primitive-mode
argument of Section~\ref{sec:quadratic-center-trace} implies
\[
 z/D\notin\mathbb Q\ \Longrightarrow\ \theta^{60}=1,
 \qquad \chi_{60}(k)=-1
 \text{ for primitive }\theta=e^{2\pi ik/60}.
\]
The units modulo sixty do not all square to one.  Consequently we do
\emph{not} conclude that individual paired twists are equal.  Instead let
$t$ be normalized relative trace to $\mathbb Q(\sqrt{15})$.  Modular
Galois symmetry replaces $\theta$ by $\theta^{u^2}$, while every square
automorphism fixes this quadratic field.  Hence
$t(\theta^{u^2})=t(\theta)$ and
$t(\theta^{2u^2})=t(\theta^2)$.  Paired row types have equal
multiplicities at each pair of relative-trace values, which is exactly
the equality used below.

\begin{lemma}\label{lem:m5-q60-hull}
For every root of unity, $(t(\theta),t(\theta^2))$ lies in the convex
hull of the following nine vectors, where $(a,b,c,d)$ denotes
$(a+bs,c+ds)$:
\begin{gather*}
 (1,0,1,0),\ (-1,0,1,0),\
 (1/2,0,-1/2,0),\ (-1/2,0,-1/2,0),\ (0,0,-1,0),\\
 (0,1/8,-1/8,0),\ (0,-1/8,-1/8,0),\
 (0,0,0,1/8),\ (0,0,0,-1/8).
\end{gather*}
For the nonrational normalized-dimension rows, the first five vertices
and $(0,-1/8,-1/8,0)$ suffice.
\end{lemma}
\begin{proof}
If the order $n$ is not divisible by sixty, the first trace is the
Ramanujan value $\mu(n)/\varphi(n)$.  For $60\mid n$, its rational
part vanishes; its irrational part is nonzero precisely when $n=60m$,
where $m$ is squarefree and coprime to thirty.  The corresponding pair is
\[
 \left(\frac{\pm s\mu(m)\chi_{60}(m)}{8\varphi(m)},
             -\frac{\mu(m)}{8\varphi(m)}\right).
\]
For $m=1$ these are the primitive-sixty vertices.  Otherwise
$w=1/\varphi(m)\leq1/6$; place weight $w$ on the matching primitive
vertex and absorb the rational second-coordinate correction in the
vertical segment $(0,[-1,1])$ contained in the rational five-vertex hull.
If only the second moment has an irrational part, $n=120m$ with the
same restrictions, and the pair lies between the two primitive-120
vertices.  All remaining pairs are rational.  Odd squarefree orders
give $(u,u)$, twice odd squarefree orders give $(-u,u)$ with
$-1/2\leq u\leq1$, and orders divisible by four give $(0,u)$ with
$-1\leq u\leq1$.  These lie in the same five-vertex hull.  The
saturated twist constraint leaves only the negative-character primitive
vertex in addition to the rational hull.  This covers all root orders;
no finite cutoff on possible twists is imposed.
\end{proof}

Let $B$ be the central induction Gram matrix in
\eqref{eq:quadratic-center-gram}.  The exact necessary integrality of
$D/z$ for a central simple of dimension $z=a+bs$ gives
\[
                  a^2-15b^2\mid60a,\qquad
                  a^2-15b^2\mid60b.
\]
The six nonunit rows in $I(1)$ have dimensions $D/f$, include the forced
row for $J$, sum to $B_{0,*}-e_0$, and leave a positive semidefinite
residual Gram matrix.  Exhaustive bounded enumeration gives $110$ frames
for $R7\_5\_443$ and $157$ for $R7\_5\_1214$.  Every bound comes
from $B$.  Before frame-specific constraints there are respectively
$3670$ and $3741$ possible residual rows.  The exact test
$vv^{\mathsf T}\leq R$ includes singular residual matrices and is also
applied to each Galois partner.

The odd-isotypic pivotal forms on $XYX^*$ and $X^3$ force all seven
second indicators to $+1$.  Combine the $21$ Gram equations with the
$12$ rational-coordinate equations for each of
\[
 \operatorname{Tr}(\theta\mid I(X_i))=0,\qquad
 \operatorname{Tr}(\theta^2\mid I(X_i))=D\nu_2(X_i),\qquad i\ne0.
\]
Multiply all columns by eight and pair Galois rows at a common hull
vertex.  The resulting system has integer columns.  A necessary integral
Gram-span condition modulo a small prime rejects respectively $13$ and
$20$ frames.  The remaining $97$ and $137$ frames are separated by
respectively $9$ and $13$ integer covectors $y$ satisfying
\[
            y\cdot c\geq0\quad\text{for every allowed column }c,
            \qquad y\cdot b<0.
\]
These inequalities rule out nonnegative real solutions, and therefore
every possible center.  The exact replay regenerates all domains and
checks the inequalities without invoking an optimizer.  Together with
Proposition~\ref{prop:m5-quadratic-galois-pu}, this proves
Theorem~\ref{thm:m5-q60}.

\subsection{Further complete signed and induction proof trees}
\label{sec:m5-complete-trees}
The same distinction between spherical obstructions and general
categorification is essential in six additional decisions.
For $R7\_5\_4530$ and $R7\_5\_4531$, every candidate spherical
dimension character either vanishes on a simple or lies in one cubic
orbit.  In that orbit a constituent of the induced unit has dimension
$D/f$.  An exact linear functional, integer-valued on the integer span
of the simple dimensions, takes the value $4/3$ on $D/f$.
Forgetting the constituent would express $D/f$ in that integer span,
a contradiction.  This dimension-order argument uses spherical
dimensions without a positivity assumption.

Complete oriented enumeration gives respectively four and one signed
quotients.  The sole quotient for $R7\_5\_4531$ is its original
tensor.  For the three nongauge quotients of $R7\_5\_4530$, the
character covers consist of a zero simple dimension or a quartic orbit
failing formal-codegree divisibility at the source codegree nine.
The gauge branches transfer sphericality by a tensor section.  Thus
both conclusions have general complex scope.

The four rings $R7\_5\_1964$, $R7\_5\_5343$, $R7\_5\_5623$
and $R7\_5\_5778$ share one induction mechanism.  Each original
spherical character cover leaves only the Frobenius--Perron Galois
orbit after zero-dimension, unequal-dual-dimension and norm obstructions.
That orbit reduces to a pseudo-unitary induction problem.  In each
ring its complete $5025$ low induction frames all fail the ordinary
induction moment inequality.  Although the certificate is handled by
the character-aware verifier, these subproofs need no character
labelling, Galois-pruning or rotation refinement.

The complete signed quotient covers are summarized below.  The last
four columns count character-component obstructions across nongauge
quotients: $Z$ is a zero simple dimension, $D$ is unequal dual
dimensions, $N$ is a norm obstruction and $F$ is a dimension-field
degree obstruction.  They are proof-tree branch counts, not numbers
of distinct fusion rings.
\begin{center}\small
\begin{tabular}{@{}l r r r r r r@{}}
\toprule
Ring & Signed quotients & Gauges & $Z$ & $D$ & $N$ & $F$\\
\midrule
$R7\_5\_1964$ & $52$ & $4$ & $160$ & $0$ & $72$ & $0$\\
$R7\_5\_5343$ & $12$ & $2$ & $22$ & $10$ & $10$ & $2$\\
$R7\_5\_5623$ & $40$ & $8$ & $96$ & $0$ & $48$ & $0$\\
$R7\_5\_5778$ & $16$ & $4$ & $24$ & $12$ & $12$ & $4$\\
\bottomrule
\end{tabular}
\end{center}
The gauge branches transfer sphericality by a tensor section and
invoke the original spherical proof.  Every remaining quotient is
covered by the displayed character obstructions.  This exhausts the
sphericalizations of a hypothetical category and proves a general
complex exclusion for each of the four rings.

The source tensors, complete proof trees, exact normal and optimized
replay receipts, and separate scope audits for this continuation are
stored under \path{evidence/rank7_m5_20261020/}.  The aggregate M5 replay
binds accepted certificates to the canonical source tensors.  Successful
finite verification establishes the stated arithmetic and enumeration;
the categorical arguments remain subject to independent mathematical
verification.
